\section{Introduction}
\label{sec:introduction}

The rapid proliferation of shared electric scooters (e-scooters) has introduced a new class of urban safety challenge. Emergency department visits related to e-scooter injuries have surged globally, with head trauma and fractures dominating presentations \citep{trivedi2019injuries, bloom2021standing, badeau2019emergency}. Fatalities, though less frequent, often involve collisions with motor vehicles at speed \citep{yang2020safety}. The power model of the speed--crash relationship predicts that even modest speed increases lead to disproportionate injury risk---a relationship especially consequential for unprotected riders \citep{elvik2013reexamination}.

Recognizing this, jurisdictions worldwide have imposed regulatory speed limits on e-scooters: 25~km/h in South Korea and France, 20~km/h in Germany and the UK, and 15--25~km/h across US cities \citep{ma2021municipal, sexton2023shared}. In practice, compliance depends almost entirely on software-based speed governors embedded in the vehicle firmware, as traditional traffic enforcement rarely extends to e-scooters. Whether these governors effectively manage speed safety is therefore a central policy question.

Yet the question ``do speed governors reduce speeding?'' is largely tautological: a device that electronically caps maximum speed will, by design, prevent riders from exceeding the governed threshold. The more consequential---and empirically open---questions concern how speed governance reshapes riding behavior on \textit{unconstrained} margins. Speed is only one dimension of riding behavior; a governor constrains it mechanically, but leaves other behavioral margins free to adjust. Drawing on the behavioral economics framework of \citet{peltzman1975effects} and subsequent risk homeostasis literature \citep{wilde1982theory, hedlund2000risky}, we ask three questions: (1)~Do governed riders compensate through more aggressive acceleration, more erratic speed profiles, or altered riding dynamics? (2)~Does governance impose a demand cost---reducing trip frequency, shortening trips, or driving platform abandonment? (3)~Does the availability of an ungoverned ``fast option'' create a behavioral escalation pathway as riders gain experience?

Prior work has been constrained in three respects. First, existing studies analyze speed behavior through a single lens---typically binary threshold exceedance or spot-speed observations \citep{cicchino2023speed, ma2021escooter}---missing the multi-dimensional nature of riding behavior and conflating mechanical speed reduction with genuinely behavioral responses. Second, no study has provided \textit{causal} evidence on how speed governance affects behavioral outcomes beyond the governed threshold itself. Third, behavioral compensation (risk homeostasis) after governor policy changes has not been tested across multiple unconstrained margins, nor has the demand-side response---whether governance reduces ridership---been examined.

We address all three gaps using a uniquely comprehensive dataset and a \textit{predict-then-validate} analytical framework. Swing, one of South Korea's largest shared e-scooter operators, offered three speed modes on identical hardware during 2023: TUB (turbo, ungoverned), STD (standard, governed to $\sim$25~km/h), and ECO (economy, governed to $\sim$20~km/h). Using 21 million GPS-instrumented trips with per-point sensor speed data from 52 cities over 10 months, we analyze eight genuinely behavioral outcomes organized into three categories: safety (harsh acceleration, harsh deceleration, speed variability), demand (trip frequency, trip distance, user retention), and riding dynamics (cruising fraction, zero-speed fraction). Speed reduction itself is treated as a mechanical manipulation check---confirmed but not claimed as a finding. Our analytical framework proceeds in two phases:

\begin{enumerate}
    \item \textbf{Phase~I---Cross-sectional prediction (February--November 2023).} The concurrent availability of governed and ungoverned modes on identical hardware provides a natural behavioral laboratory. Using generalized estimating equations (GEE), within-subject paired comparisons (76,726 mode-switcher pairs), and SHAP-based feature importance from gradient-boosted tree models, we characterize which behavioral margins are strongly mode-dependent and which are not. From these cross-sectional patterns, we derive six testable predictions about what removing the ungoverned mode would do to each outcome---specifying both the expected direction and whether the effect should be significant or null.

    \item \textbf{Phase~II---Causal validation (December 2023 natural experiment).} In December 2023, Swing banned the ungoverned TUB mode system-wide, providing an exogenous policy shock. We validate each prediction using multi-outcome difference-in-differences with continuous treatment intensity and Bonferroni correction, a within-user mode-switcher compensation test, and dose-response analysis across 50 cities. Survival analysis with multiple event types quantifies how ungoverned mode availability accelerates behavioral escalation over time.
\end{enumerate}

We formalize six testable predictions from the observational data and evaluate each against the causal evidence (Table~\ref{tab:predictions} and Table~\ref{tab:validation}). All six predictions are confirmed: cross-sectional behavioral profiling accurately forecasts which margins respond to governance removal and which do not.

To our knowledge, this is the first study to demonstrate that cross-sectional behavioral profiling \textit{predicts} causal policy outcomes for speed governance---beyond the tautological speed reduction. The predict-then-validate framework establishes that multi-outcome observational analysis has genuine forecasting power for policy effects, providing a methodological template for pre-implementation assessment of speed management interventions.
